\documentclass[a4paper,12pt]{article}

% -------------------------------------------------
% Кодировка, язык, математика
% -------------------------------------------------
\usepackage[T2A]{fontenc}
\usepackage[utf8]{inputenc}
\usepackage[russian]{babel}
\usepackage{amsmath,amssymb,amsthm,mathtools}
\usepackage{helvet}
\renewcommand{\familydefault}{\sfdefault}
\usepackage{icomma}

% -------------------------------------------------
% Вёрстка и типографика
% -------------------------------------------------
\usepackage[a4paper,margin=2.15cm]{geometry}
\usepackage{microtype}
\usepackage{parskip}
\usepackage{setspace}
\onehalfspacing
\usepackage{enumitem}
\usepackage{tabularx,booktabs,longtable}
\usepackage{xcolor}
\usepackage{fancyhdr}
\usepackage{hyperref}

% -------------------------------------------------
% Графика
% -------------------------------------------------
\usepackage{tikz}
\usetikzlibrary{calc}
\usepackage{tkz-euclide}
\usepackage{pgfplots}
\pgfplotsset{compat=1.18}

% -------------------------------------------------
% Цветовая схема
% -------------------------------------------------
\definecolor{accent}{HTML}{008080}
\definecolor{darkaccent}{HTML}{004D4D}
\definecolor{textgray}{HTML}{333333}
\definecolor{softgray}{HTML}{707070}
\definecolor{linegray}{HTML}{D7E1E1}

% -------------------------------------------------
% Общие настройки
% -------------------------------------------------
\color{textgray}
\setlength{\parindent}{0pt}
\setlength{\parskip}{0.55em}
\setlist[itemize]{leftmargin=1.5em,itemsep=0.25em,topsep=0.25em}
\setlist[enumerate]{leftmargin=1.8em,itemsep=0.45em,topsep=0.35em}

\hypersetup{
    colorlinks=true,
    linkcolor=darkaccent,
    urlcolor=darkaccent,
    pdftitle={Чек-лист по производной},
    pdfauthor={Лёвин Артём Александрович}
}

\pagestyle{fancy}
\fancyhf{}
\fancyhead[L]{\small\color{softgray}Производная функции}
\fancyhead[R]{\small\color{softgray}ЕГЭ}
\fancyfoot[C]{\small\color{softgray}\thepage}
\renewcommand{\headrulewidth}{0.4pt}
\renewcommand{\headrule}{\hbox to\headwidth{\color{linegray}\leaders\hrule height \headrulewidth\hfill}}

% -------------------------------------------------
% Служебные команды
% -------------------------------------------------
\newcommand{\sectionline}{%
    \vspace{0.25em}
    {\color{accent}\rule{\linewidth}{0.8pt}}
    \vspace{0.45em}
}

\newcommand{\keyformula}[1]{%
    \[
    \boxed{\displaystyle #1}
    \]
}

\newcommand{\checkitem}[1]{%
    \par\noindent
    {\color{accent}\large$\square$}\hspace{0.55em}#1
}

\begin{document}

% =================================================
% ТИТУЛЬНЫЙ ЛИСТ
% =================================================
\begin{titlepage}
\thispagestyle{empty}

\begin{center}

\vspace*{2.2cm}

{\color{darkaccent}\Large\bfseries ЧЕК-ЛИСТ}

\vspace{0.8cm}

{\Huge\bfseries Производная функции}

\vspace{0.45cm}

{\Large
Основные формулы, произведение и частное,\\
тригонометрические, показательные\\
и логарифмические функции
}

\vspace{1.6cm}

{\color{accent}\rule{0.62\linewidth}{1pt}}

\vspace{1.6cm}

{\large Подготовка к ЕГЭ по профильной математике}

\vfill

{\large
Лёвин Артём Александрович\\[0.35em]
эксклюзивно для Николь Саркисьянц
}

\vspace{0.65cm}

{\large \today}

\vspace{1.8cm}

\end{center}

% Сдержанная кривая Безье
\begin{tikzpicture}[remember picture,overlay]
\draw[
    accent,
    line width=1.2pt
]
($(current page.south west)+(1.6cm,1.45cm)$)
.. controls
($(current page.south west)+(5.3cm,2.15cm)$)
and
($(current page.south east)+(-5.2cm,0.75cm)$)
..
($(current page.south east)+(-1.6cm,1.45cm)$);
\end{tikzpicture}

\end{titlepage}

% =================================================
% 1. БАЗОВЫЕ ФОРМУЛЫ
% =================================================
\section*{\color{darkaccent}1. Базовые производные}
\sectionline

Производная --- функция, характеризующая мгновенную скорость изменения исходной функции. 
В вычислительных задачах основой служит таблица производных и правила дифференцирования.

\subsection*{\color{accent}Степенная функция}

Для постоянного показателя степени \(n\):

\keyformula{(x^n)'=n x^{\,n-1}}

Формула работает также для отрицательных и дробных степеней там, где соответствующая функция определена.

\[
(x^{-10})'=-10x^{-11}.
\]

\[
(\sqrt{x})'
=
\left(x^{1/2}\right)'
=
\frac{1}{2}x^{-1/2}
=
\frac{1}{2\sqrt{x}},
\qquad x>0.
\]

\[
(x\sqrt{x})'
=
\left(x^{3/2}\right)'
=
\frac32\sqrt{x}.
\]

\subsection*{\color{accent}Постоянная}

Если \(C\) --- число, то

\keyformula{C'=0}

В частности, если параметр \(p\) не зависит от \(x\), то

\[
(p^2)'=0.
\]

При этом

\[
(p^2x^{10})'
=
p^2\cdot 10x^9.
\]

Параметр \(p\) при дифференцировании по \(x\) рассматривается как постоянный множитель.

\subsection*{\color{accent}Сумма и разность}

\keyformula{(u\pm v)'=u'\pm v'}

Например,

\[
(p^2+x^{10})'=0+10x^9=10x^9.
\]

\subsection*{\color{accent}Постоянный множитель}

\keyformula{(C f(x))'=C f'(x)}

Например,

\[
(3x^{10})'=30x^9.
\]

\checkitem{Умею отличать переменную \(x\) от постоянного параметра.}

\checkitem{Помню, что при дифференцировании степень уменьшается на \(1\).}

\checkitem{Помню, что производная любой постоянной равна нулю.}

% =================================================
% 2. ПРОИЗВЕДЕНИЕ
% =================================================
\section*{\color{darkaccent}2. Производная произведения}
\sectionline

Если две функции зависят от \(x\) и перемножаются, используется формула произведения:

\keyformula{(uv)'=u'v+uv'}

Удобный алгоритм:

\begin{enumerate}
    \item Выделить первый множитель \(u\) и второй множитель \(v\).
    \item Найти \(u'\) и \(v'\).
    \item Подставить в формулу \(u'v+uv'\).
    \item При необходимости упростить результат.
\end{enumerate}

\subsection*{\color{accent}Пример 1}

\[
y=x^3\cdot x^{11}.
\]

Пусть

\[
u=x^3,\qquad v=x^{11}.
\]

Тогда

\[
u'=3x^2,\qquad v'=11x^{10}.
\]

Следовательно,

\[
y'
=
3x^2\cdot x^{11}
+
x^3\cdot 11x^{10}
=
14x^{13}.
\]

Для одночленов тот же ответ можно получить предварительным упрощением:

\[
x^3\cdot x^{11}=x^{14},
\qquad
(x^{14})'=14x^{13}.
\]

\subsection*{\color{accent}Пример 2}

\[
y=x^3(5x-2).
\]

Берём

\[
u=x^3,\qquad v=5x-2.
\]

Тогда

\[
u'=3x^2,\qquad v'=5.
\]

По формуле произведения:

\[
y'
=
3x^2(5x-2)+5x^3.
\]

Можно упростить:

\[
y'=20x^3-6x^2.
\]

\subsection*{\color{accent}Как распознать формулу}

Сигнал к применению правила произведения --- наличие двух выражений, каждое из которых зависит от \(x\):

\[
x^2\sin x,
\qquad
(2x-1)(x^5+3x^{10}),
\qquad
x^2e^{3x}.
\]

\checkitem{При произведении дифференцирую по очереди каждый множитель.}

\checkitem{Между двумя полученными произведениями ставлю знак \(+\).}

% =================================================
% 3. ЧАСТНОЕ
% =================================================
\section*{\color{darkaccent}3. Производная частного}
\sectionline

Если

\[
y=\frac{u}{v},
\qquad v\ne0,
\]

то

\keyformula{\left(\frac{u}{v}\right)'=\frac{u'v-uv'}{v^2}}

Порядок в числителе принципиален:

\[
u'v-uv'.
\]

Менять эти два слагаемых местами нельзя.

\subsection*{\color{accent}Пример}

\[
y=\frac{x^4}{x^{12}},
\qquad x\ne0.
\]

Пусть

\[
u=x^4,\qquad v=x^{12}.
\]

Тогда

\[
u'=4x^3,
\qquad
v'=12x^{11}.
\]

Следовательно,

\[
y'
=
\frac{4x^3\cdot x^{12}-x^4\cdot12x^{11}}
{(x^{12})^2}.
\]

То есть

\[
y'
=
\frac{-8x^{15}}{x^{24}}
=
-\frac{8}{x^9},
\qquad x\ne0.
\]

То же можно проверить после упрощения исходной функции:

\[
\frac{x^4}{x^{12}}=x^{-8},
\qquad
(x^{-8})'=-8x^{-9}.
\]

\subsection*{\color{accent}Область допустимых значений}

Если в исходной функции есть знаменатель, необходимо сохранить условие

\[
v(x)\ne0.
\]

Например, для

\[
\frac{x^3}{5x-1}
\]

имеем

\[
5x-1\ne0
\quad\Longrightarrow\quad
x\ne\frac15.
\]

\checkitem{В формуле частного помню порядок \(u'v-uv'\).}

\checkitem{В знаменателе ставлю квадрат исходного \(v\), а не квадрат \(v'\).}

\checkitem{До преобразований фиксирую ограничения знаменателя.}

% =================================================
% 4. ТРИГОНОМЕТРИЧЕСКИЕ ФУНКЦИИ
% =================================================
\section*{\color{darkaccent}4. Производные тригонометрических функций}
\sectionline

Основные формулы:

\[
\begin{aligned}
(\sin x)'&=\cos x,\\
(\cos x)'&=-\sin x,\\
(\tg x)'&=\frac{1}{\cos^2x},\\
(\ctg x)'&=-\frac{1}{\sin^2x}.
\end{aligned}
\]

\subsection*{\color{accent}Почему \((\tg x)'=\dfrac{1}{\cos^2x}\)}

Используем

\[
\tg x=\frac{\sin x}{\cos x},
\qquad
\cos x\ne0.
\]

По формуле частного:

\[
(\tg x)'
=
\frac{\cos x\cdot\cos x-\sin x(-\sin x)}
{\cos^2x}.
\]

Значит,

\[
(\tg x)'
=
\frac{\cos^2x+\sin^2x}{\cos^2x}.
\]

С учётом тождества

\[
\sin^2x+\cos^2x=1
\]

получаем

\keyformula{(\tg x)'=\frac{1}{\cos^2x}}

\subsection*{\color{accent}Почему \((\ctg x)'=-\dfrac{1}{\sin^2x}\)}

\[
\ctg x=\frac{\cos x}{\sin x},
\qquad
\sin x\ne0.
\]

Тогда

\[
(\ctg x)'
=
\frac{(-\sin x)\sin x-\cos x\cos x}
{\sin^2x}.
\]

Следовательно,

\[
(\ctg x)'
=
-\frac{\sin^2x+\cos^2x}{\sin^2x}
=
-\frac{1}{\sin^2x}.
\]

\subsection*{\color{accent}Пример произведения}

\[
y=x\sin x.
\]

Берём

\[
u=x,
\qquad
v=\sin x.
\]

Тогда

\[
u'=1,
\qquad
v'=\cos x.
\]

Следовательно,

\[
y'=\sin x+x\cos x.
\]

\checkitem{Помню знак минус в производной \(\cos x\).}

\checkitem{Помню знак минус в производной \(\ctg x\).}

\checkitem{Умею выводить формулы для \(\tg x\) и \(\ctg x\) через правило частного.}

% =================================================
% 5. СЛОЖНАЯ ФУНКЦИЯ
% =================================================
\section*{\color{darkaccent}5. Производная сложной функции}
\sectionline

Во многих формулах занятия возникает одна общая идея: необходимо продифференцировать внешнюю функцию и умножить результат на производную внутренней функции.

Если

\[
y=F(g(x)),
\]

то

\keyformula{y'=F'(g(x))\cdot g'(x)}

Это правило часто называют правилом цепочки.

Например,

\[
y=e^{x^2-x-6}.
\]

Внутренняя функция:

\[
g(x)=x^2-x-6,
\qquad
g'(x)=2x-1.
\]

Следовательно,

\[
y'=(2x-1)e^{x^2-x-6}.
\]

\checkitem{При наличии функции ``внутри функции'' ищу внутреннее выражение.}

\checkitem{После дифференцирования внешней функции умножаю на производную внутренней.}

% =================================================
% 6. ПОКАЗАТЕЛЬНАЯ ФУНКЦИЯ
% =================================================
\section*{\color{darkaccent}6. Показательная функция}
\sectionline

Для основания

\[
a>0,\qquad a\ne1
\]

справедлива формула

\keyformula{\bigl(a^{f(x)}\bigr)'=f'(x)a^{f(x)}\ln a}

\subsection*{\color{accent}Пример}

\[
y=2^{x^2-x-6}.
\]

Производная показателя:

\[
(x^2-x-6)'=2x-1.
\]

Поэтому

\[
y'
=
(2x-1)\,2^{x^2-x-6}\ln2.
\]

\subsection*{\color{accent}Экспонента}

При основании \(e\):

\[
\ln e=1.
\]

Поэтому формула упрощается:

\keyformula{\bigl(e^{f(x)}\bigr)'=f'(x)e^{f(x)}}

Например,

\[
\left(e^{x^2-x-6}\right)'
=
(2x-1)e^{x^2-x-6}.
\]

\subsection*{\color{accent}Комбинация с произведением}

\[
y=x^2e^{3x}.
\]

Пусть

\[
u=x^2,
\qquad
v=e^{3x}.
\]

Тогда

\[
u'=2x,
\]

а по правилу цепочки

\[
v'=3e^{3x}.
\]

Следовательно,

\[
y'
=
2xe^{3x}+3x^2e^{3x}.
\]

Можно вынести общий множитель:

\[
y'=xe^{3x}(2+3x).
\]

\checkitem{В \(a^{f(x)}\) умножаю результат на \(\ln a\).}

\checkitem{Для основания \(e\) множитель \(\ln e\) равен \(1\).}

% =================================================
% 7. ЛОГАРИФМИЧЕСКАЯ ФУНКЦИЯ
% =================================================
\section*{\color{darkaccent}7. Логарифмическая функция}
\sectionline

Для

\[
a>0,\qquad a\ne1,\qquad f(x)>0
\]

имеем

\keyformula{\bigl(\log_a f(x)\bigr)'=
\frac{f'(x)}{f(x)\ln a}}

Ограничение

\[
f(x)>0
\]

является частью ОДЗ исходной логарифмической функции.

\subsection*{\color{accent}Пример 1}

\[
y=\log_2(x^2-x-6).
\]

ОДЗ:

\[
x^2-x-6>0.
\]

Разложим:

\[
(x-3)(x+2)>0,
\]

откуда

\[
x\in(-\infty,-2)\cup(3,+\infty).
\]

Производная:

\[
y'
=
\frac{2x-1}
{(x^2-x-6)\ln2}.
\]

\subsection*{\color{accent}Пример 2}

\[
y=\log_5(4x+3).
\]

ОДЗ:

\[
4x+3>0
\quad\Longrightarrow\quad
x>-\frac34.
\]

Производная:

\[
y'
=
\frac{4}{(4x+3)\ln5}.
\]

\subsection*{\color{accent}Десятичный логарифм}

\[
\lg x=\log_{10}x.
\]

Поэтому

\keyformula{(\lg f(x))'=
\frac{f'(x)}{f(x)\ln10}}

Например,

\[
(\lg(4x+3))'
=
\frac{4}{(4x+3)\ln10}.
\]

\subsection*{\color{accent}Натуральный логарифм}

\[
\ln x=\log_e x.
\]

Так как \(\ln e=1\),

\keyformula{(\ln f(x))'=\frac{f'(x)}{f(x)}}

Например,

\[
y=\ln(5x^2).
\]

ОДЗ:

\[
5x^2>0
\quad\Longrightarrow\quad
x\ne0.
\]

Тогда

\[
y'
=
\frac{10x}{5x^2}
=
\frac{2}{x},
\qquad x\ne0.
\]

\checkitem{Перед дифференцированием логарифма фиксирую условие \(f(x)>0\).}

\checkitem{В знаменателе формулы для \(\log_a f(x)\) присутствует \(\ln a\).}

\checkitem{Для натурального логарифма множитель \(\ln e\) исчезает.}

% =================================================
% 8. СТРАТЕГИЯ РАСПОЗНАВАНИЯ
% =================================================
\section*{\color{darkaccent}8. Как выбрать формулу}
\sectionline

Перед вычислением производной полезно сначала определить структуру функции.

\begin{center}
\renewcommand{\arraystretch}{1.4}
\begin{tabularx}{\linewidth}{@{}>{\raggedright\arraybackslash}p{0.34\linewidth}X@{}}
\toprule
\textbf{Структура} & \textbf{Главное правило} \\
\midrule
\(u(x)+v(x)\) & дифференцировать каждое слагаемое отдельно \\
\(C\,u(x)\) & постоянный множитель оставить перед производной \\
\(u(x)v(x)\) & использовать \((uv)'=u'v+uv'\) \\
\(\dfrac{u(x)}{v(x)}\) & использовать правило частного и проверить \(v(x)\ne0\) \\
\(F(g(x))\) & применить правило цепочки \\
\(a^{f(x)}\) & производная показателя, исходная степень, \(\ln a\) \\
\(\log_a f(x)\) & \(\dfrac{f'(x)}{f(x)\ln a}\), обязательно \(f(x)>0\) \\
\bottomrule
\end{tabularx}
\end{center}

\subsection*{\color{accent}Главный алгоритм}

\begin{enumerate}
    \item Посмотрите на функцию целиком.
    \item Найдите главное действие: сумма, произведение, деление, степень, логарифм.
    \item Выберите соответствующую формулу.
    \item Разбейте сложное выражение на простые части.
    \item Найдите производные этих частей.
    \item Соберите результат по выбранной формуле.
    \item Упростите ответ.
    \item Проверьте ОДЗ исходной функции.
\end{enumerate}

% =================================================
% 9. ТИПИЧНЫЕ ОШИБКИ
% =================================================
\section*{\color{darkaccent}9. Что особенно важно контролировать}
\sectionline

\begin{enumerate}
    \item
    \textbf{Степень.}
    \[
    (x^n)'=nx^{n-1}.
    \]
    Например,
    \[
    (x^{-10})'=-10x^{-11}.
    \]

    \item
    \textbf{Параметры.}
    Если \(p\) --- постоянная относительно \(x\), то
    \[
    (p^2)'=0.
    \]

    \item
    \textbf{Произведение.}
    Нельзя писать
    \[
    (uv)'=u'v'.
    \]
    Правильно:
    \[
    (uv)'=u'v+uv'.
    \]

    \item
    \textbf{Частное.}
    В числителе сохраняется порядок
    \[
    u'v-uv'.
    \]

    \item
    \textbf{Косинус.}
    \[
    (\cos x)'=-\sin x.
    \]

    \item
    \textbf{Логарифм.}
    Аргумент логарифма должен быть положительным:
    \[
    f(x)>0.
    \]

    \item
    \textbf{Знаменатель.}
    Любой знаменатель должен быть отличен от нуля.

    \item
    \textbf{Скобки.}
    Если квадрат относится ко всему выражению, это должно быть видно:
    \[
    (\ln x)^2,
    \qquad
    (\cos x)^2.
    \]

    \item
    \textbf{Запись дроби.}
    Если множитель относится ко всей дроби, запись должна быть однозначной:
    \[
    x^2\cdot\frac{1}{2\sqrt{x}}
    =
    \frac{x^2}{2\sqrt{x}}.
    \]

    \item
    \textbf{Формулы полезно узнавать в обе стороны.}
    Например,
    \[
    \sin^2x+\cos^2x=1
    \]
    и
    \[
    1-\sin^2x=\cos^2x.
    \]
\end{enumerate}

% =================================================
% 10. КРАТКАЯ ТАБЛИЦА
% =================================================
\section*{\color{darkaccent}10. Формулы, которые нужно знать наизусть}
\sectionline

\begin{center}
\renewcommand{\arraystretch}{1.5}
\begin{tabularx}{\linewidth}{@{}X X@{}}
\toprule
\textbf{Функция} & \textbf{Производная} \\
\midrule
\(C\) & \(0\) \\

\(x^n\) & \(nx^{n-1}\) \\

\(\sqrt{x}\) & \(\dfrac{1}{2\sqrt{x}}\) \\

\(\sin x\) & \(\cos x\) \\

\(\cos x\) & \(-\sin x\) \\

\(\tg x\) & \(\dfrac{1}{\cos^2x}\) \\

\(\ctg x\) & \(-\dfrac{1}{\sin^2x}\) \\

\(uv\) & \(u'v+uv'\) \\

\(\dfrac{u}{v}\) & \(\dfrac{u'v-uv'}{v^2}\) \\

\(a^{f(x)}\) & \(f'(x)a^{f(x)}\ln a\) \\

\(e^{f(x)}\) & \(f'(x)e^{f(x)}\) \\

\(\log_a f(x)\) & \(\dfrac{f'(x)}{f(x)\ln a}\) \\

\(\ln f(x)\) & \(\dfrac{f'(x)}{f(x)}\) \\
\bottomrule
\end{tabularx}
\end{center}

% =================================================
% ТРЕНИРОВОЧНЫЙ БЛОК
% =================================================
\clearpage
\section*{\color{darkaccent}Тренировочный блок}
\sectionline

Найдите производную каждой функции. Где это требуется, сначала укажите область допустимых значений исходной функции.

\subsection*{\color{accent}Средний уровень}

\begin{enumerate}[label=\textbf{\arabic*)}]
    \item
    \[
    y=x^4(3x-2).
    \]

    \item
    \[
    y=\frac{x^3}{4x-1}.
    \]

    \item
    \[
    y=x^2\sin x.
    \]
\end{enumerate}

\subsection*{\color{accent}Выше среднего}

\begin{enumerate}[label=\textbf{\arabic*)},resume]
    \item
    \[
    y=(2x-1)(x^5+3x^2).
    \]

    \item
    \[
    y=\frac{x^4+1}{x^2}.
    \]

    \item
    \[
    y=x^2e^{3x}.
    \]

    \item
    \[
    y=5^{\,2x^2-x}.
    \]

    \item
    \[
    y=\log_3(x^2-4x+3).
    \]

    \item
    \[
    y=\sin x\cdot\cos x.
    \]
\end{enumerate}

\subsection*{\color{accent}Сложный уровень}

\begin{enumerate}[label=\textbf{\arabic*)},resume]
    \item
    \[
    y=\frac{x^2e^{2x}}{\ln x}.
    \]
\end{enumerate}

% =================================================
% ОТВЕТЫ
% =================================================
\clearpage
\section*{\color{darkaccent}Ответы к тренировочному блоку}
\sectionline

\begin{center}
\renewcommand{\arraystretch}{1.55}
\begin{tabularx}{\linewidth}{@{}p{0.07\linewidth}X@{}}
\toprule
\textbf{№} & \textbf{Ответ} \\
\midrule

1 &
\[
y'=4x^3(3x-2)+3x^4
=15x^4-8x^3.
\]
\\

2 &
ОДЗ:
\[
x\ne\frac14.
\]
\[
y'
=
\frac{3x^2(4x-1)-4x^3}{(4x-1)^2}
=
\frac{x^2(8x-3)}{(4x-1)^2}.
\]
\\

3 &
\[
y'=2x\sin x+x^2\cos x.
\]
\\

4 &
\[
y'
=
2(x^5+3x^2)
+
(2x-1)(5x^4+6x).
\]
\\

5 &
ОДЗ:
\[
x\ne0.
\]
\[
y'
=
\frac{4x^3x^2-(x^4+1)2x}{x^4}
=
2x-\frac{2}{x^3}.
\]
\\

6 &
\[
y'
=
2xe^{3x}+3x^2e^{3x}
=
xe^{3x}(2+3x).
\]
\\

7 &
\[
y'
=
(4x-1)5^{\,2x^2-x}\ln5.
\]
\\

8 &
ОДЗ:
\[
x^2-4x+3>0.
\]
То есть
\[
(x-1)(x-3)>0,
\]
поэтому
\[
x\in(-\infty,1)\cup(3,+\infty).
\]
Производная:
\[
y'
=
\frac{2x-4}
{(x^2-4x+3)\ln3}.
\]
\\

9 &
\[
y'
=
\cos^2x-\sin^2x
=
\cos2x.
\]
\\

10 &
ОДЗ:
\[
x>0,\qquad \ln x\ne0,
\]
то есть
\[
x>0,\qquad x\ne1.
\]

Пусть
\[
u=x^2e^{2x},
\qquad
v=\ln x.
\]

Тогда
\[
u'
=
2xe^{2x}+2x^2e^{2x}
=
2xe^{2x}(1+x),
\]
\[
v'=\frac1x.
\]

Следовательно,
\[
y'
=
\frac{
2xe^{2x}(1+x)\ln x
-
x^2e^{2x}\cdot\frac1x
}
{(\ln x)^2}.
\]

После упрощения:
\[
\boxed{
y'
=
\frac{
xe^{2x}\bigl(2(1+x)\ln x-1\bigr)
}
{(\ln x)^2}
}.
\]
\\

\bottomrule
\end{tabularx}
\end{center}

\end{document}